\chapter*{如何阅读本书}

如果你刚开始学 \cpp，可以按顺序读第一到第七章。不要急着跳到模板、并发和性能；如果你还不能解释一个对象何时构造、何时析构、谁拥有资源，那么高级语法只会放大错误。

如果你已经会写基本程序，可以重点读第五章到第十三章。读的时候建议做三件事：第一，手敲每章的核心代码；第二，故意改坏一个边界条件，看编译器、测试或运行时如何反馈；第三，把本章的接口改成自己的小工具，而不是只复制书里的函数。

每章会反复出现几个固定栏目：

\begin{itemize}
  \item \textbf{问题入口}：先给一个真实的小需求。
  \item \textbf{朴素写法}：先写能想到的直接方案。
  \item \textbf{失败原因}：解释它为什么会错、慢或难维护。
  \item \textbf{现代写法}：用标准库和语言机制修正。
  \item \textbf{边界检查}：列出容易遗漏的输入、生命周期和复杂度。
\end{itemize}

\begin{warningbox}
不要把书中代码当成唯一答案。真正要学的是代码背后的不变量：哪些数据必须一直成立，哪些资源只有一个拥有者，哪些函数只观察不修改，哪些错误必须显式返回或抛出。
\end{warningbox}
